\begin{table}[!htbp]
\centering
\caption{Synthetic-DiD robustness excluding October 2022 from the pre-treatment window}
\label{tab:v1_sdid_drop_last_pre}
\begin{threeparttable}
\small
\begin{tabular}{lc}
\toprule
& \# of registered unemployed \
\midrule
ATT: adjustment period & 0.037 \
Impact (percent): adjustment period & 3.8 \
ATT: later period & 0.056 \
Impact (percent): later period & 5.7 \
Largest month weight & 0.651 \
Month with largest weight & Sep 2022 (event time -2) \
Share of weights in event times -6 to -2 & 0.918 \
\bottomrule
\end{tabular}
\begin{tablenotes}[flushleft]
\footnotesize
\item \emph{Notes:} This robustness re-estimates the adjusted expanded-donor synthetic-DiD after removing October 2022 (event time -1) from the pre-treatment window, then rebalancing the panel over the remaining months. The phase-specific unemployment effects are reported as point estimates only because the robustness is run with \texttt{vce(noinference)} while keeping projected CNO1-by-month covariates. The concentration statistic sums weights over event times -6 through -2 because event time -1 is excluded by construction.
\end{tablenotes}
\end{threeparttable}
\end{table}
